% !TeX encoding = UTF-8
\chapter{Conclusion}

In the context of this thesis, a multi-component framework for assessing the quality of an LLM-based medical chatbot in a prenatal care consultation setting was designed, implemented, and evaluated. The project addressed four distinct evaluation targets: information extraction quality, topic coverage, citation faithfulness, and mandatory question compliance, each approached with multiple complementary methods to ensure robust and reproducible findings.

\section{Summary of Contributions}

The first contribution of this thesis is a systematic benchmarking of four information extraction strategies across four Large Language Models using two evaluation approaches: thresholded semantic matching and the adapted SUSWIR metric. 
The results of the evaluation demonstrate that prompt complexity and model capability interact in non-trivial ways: larger frontier models perform best with minimal prompting, while smaller open-source models benefit from structured extraction frameworks. 
This finding has direct practical implications for teams designing cost-efficient medical NLP pipelines.

The second contribution is a comprehensive evaluation of topic coverage across two synthetic conversation datasets using three independent approaches: Hungarian matching, bi-encoder with entailment, and LLM-as-a-Judge. 
The consistent agreement across methods, confirmed by statistically significant Spearman rank correlations, validates the multi-method evaluation design and provides a reliable picture of the chatbot's information delivery behavior.

The third and most clinically significant contribution is the demonstration that a supervised dialogue management layer substantially improves mandatory question compliance and topic coverage without degrading citation faithfulness. 
The supervised architecture raises acceptable compliance to 100\% in the smaller dataset and produces a 21-percentage-point recall gain in the larger dataset, providing strong empirical support for the necessity of active conversational scaffolding in safety-critical medical applications.

\section{Research Objectives Revisited}

The first objective \textbf{comparing existing information extraction techniques against medical texts and identifying the best-performing approach} was addressed in Section 5.1. 
A systematic benchmarking of four extraction strategies across four language models identified GPT-5 Mini with Naive prompting as the most suitable configuration, balancing F1 score with intrinsic quality as measured by the adapted SUSWIR metric. 
The open-source alternative Ministral-14B with UIE achieved marginally higher F1 and is recommended when maximizing recall alignment to a fixed ground truth is the primary objective.

The second objective \textbf{measuring the coverage of medically relevant information provided by the chatbot relative to the required ground-truth topics} was addressed in Section 5.2. 
The LLM Judge consistently reported the highest coverage estimates, with supervised dialogue management improving hit rates and weighted critical recall across all datasets and patient personas.

The third objective \textbf{evaluating citation accuracy} was addressed in Section 5.3.
Strict citation precision of 56--62\% and relaxed precision of 69--78\% were observed across datasets, with supervised dialogue management marginally improving faithfulness and consistently reducing the proportion of unsupported claims.

The fourth objective \textbf{evaluating the dialogue state management ability of the chatbot} was addressed in Section 5.4. 
The supervised manager substantially outperformed the naive configuration on mandatory question compliance, raising acceptable compliance to 100\% in Dataset A and increasing mean question recall by 21 percentage points in Dataset B, while also improving topic coverage across all patient personas.

\section{Impact and Outlook}

The findings of this thesis demonstrate that general-purpose LLMs, when augmented with domain-specific retrieval, structured knowledge extraction, and active dialogue supervision, can serve as a foundation for safe and comprehensive medical consultation systems. 
While the results are grounded in synthetic evaluation and require clinical validation before deployment, they provide a concrete methodological blueprint for researchers and practitioners working at the intersection of natural language processing and healthcare AI.
The proposed framework is designed to be modular and extensible, supporting future adaptation to new clinical domains, languages, and patient populations.